\section{Physical realization of anchored finite contact-jet families}
\label{sec:v22-anchored-realization}

The statistical experiment below is formulated in one fixed laboratory chart.
We record that the resulting finite-dimensional contact-jet models are
realized by actual nearby billiard tables.  The contact patch used by the
statistical theorem is chosen inside the region where every cutoff is exactly
one, so the displayed polynomial jet family is literally the physical graph
on the whole observation patch.

\begin{proposition}[Anchored realization of finite labelled jets]
\label{prop:v22-anchored-realization}
Let a smooth periodic dispersing billiard contain a selected labelled closest
pair with positive gap $g_0$, positive contact curvatures, and a strict
clearance margin separating this pair from every competing collision segment
outside fixed contact neighborhoods.  Fix $M\ge2$.  Choose one Euclidean
laboratory chart in which the reference contacts have common transverse
coordinate $y=0$ and common tangent direction.

There exist numbers $0<u_*<u_{**}$ and a smooth finite-dimensional family of
actual periodic billiard tables through the reference table, with coordinates
\[
 (g,q_{0,2},\ldots,q_{0,M},q_{1,2},\ldots,q_{1,M}),
\]
such that on the whole statistical patch $|y|<u_*$ the two participating
boundaries are exactly
\begin{equation}\label{eq:v22-realized-anchored-graphs}
 (-\psi_0(u),u),\qquad (g+\psi_1(v),v),
\end{equation}
with
\[
 \psi_b(0)=\psi_b'(0)=0,
 \qquad \psi_b^{(n)}(0)=q_{b,n}\quad(2\le n\le M).
\]
The labels, transverse origin, and tangent orientation are the same throughout
the family.  Strict convexity, embeddedness of every participating obstacle,
disjointness of distinct obstacle closures and of their relevant lattice
translates, and the selected-channel clearance margin persist uniformly after
the parameter neighborhood is made sufficiently small.

The same statement holds if the two contacts lie on different labelled
obstacles, on two lattice translates, or in two disjoint contact
neighborhoods of one connected periodic obstacle, provided lattice-related
perturbations are made equivariantly.
\end{proposition}

\begin{proof}
Choose disjoint arclength neighborhoods of the two reference contacts so
small that the fixed laboratory graph representation is valid.  Select smooth
cutoffs $\chi_b$ supported in $|y|<u_{**}$ and identically one on a smaller
interval $|y|<2u_*$.  We choose the statistical endpoint box only afterwards,
inside $|y|<u_*$.  Thus every graph used by the physical statistical theorem
lies strictly inside the plateau of the cutoff.

For small parameters set
\begin{equation}\label{eq:v22-local-bump}
 \delta\psi_b(y)
 =\chi_b(y)\sum_{n=2}^M\frac{\delta q_{b,n}}{n!}y^n.
\end{equation}
On $|y|<u_*$ this is exactly the displayed polynomial perturbation and hence
changes precisely the prescribed finite jet without moving the contact
origin or tangent line.  A gap variation is realized by a normal displacement
of the second patch which is constant on $|y|<2u_*$ and is cut off before the
boundary leaves the chosen coordinate neighborhood.  Its constant value on
the statistical patch is absorbed into $g$ in
\eqref{eq:v22-realized-anchored-graphs}.

If the two contacts lie on one connected obstacle, take the two cutoff
supports disjoint.  The remaining part of the boundary is unchanged.  A
sufficiently small $C^1$ perturbation of an embedded compact curve remains
embedded; a sufficiently small $C^2$ perturbation preserves strict positive
curvature.  The reference table has positive separation of distinct obstacle
closures and a strict channel-isolation margin away from the chosen contact
patches.  On a compact fundamental cell these inequalities persist under
small $C^1$ perturbations.  For periodic translates, the unchanged lattice
and a uniform obstacle diameter leave only finitely many nearby translates to
check.  Shrinking the parameter ball therefore preserves global embeddedness,
disjointness and channel isolation, including when both patches belong to the
same connected obstacle.

Because each cutoff is constant on the statistical patch, the map from bump
coefficients to the displayed contact jets has identity Jacobian.  The
inverse function theorem gives an open finite-dimensional coordinate family.
\end{proof}

\begin{remark}[Analytic versus smooth realization]
The proposition serves the finite-dimensional statistical local experiment,
which needs only finite smoothness on compact parameter sets.  The analytic
continuation conclusions of Theorems~\ref{thm:v22-signed-rigidity} and
\ref{thm:v23-intrinsic-table-rigidity} are separate deterministic
statements for analytic boundary germs; no compactly supported analytic bump
is used there.
\end{remark}
